% EN49 c09 | 原书 p.87–97（PDF 98–108） | 第 6 章 §6.1–6.2
% 骨架文件：标题行已设定，正文由译者填写。每页起首用 % ===== 原书 p.N ===== 注释标记。
\section{分次}

\subsection{矩阵分解范畴}

% ===== 原书 p.87 =====
设 $X$ 是光滑复拟射影簇，$W\colon X\to\C$ 是正则函数，在 $0$ 上只有一条临界纤维
$X_0$。与 LG 模型 $(X,W)$ 自然相关的有两个范畴：

\begin{definition}\label{def:sing-cat}
奇点范畴 $\operatorname{Sing}(X_0)$，它是 $\operatorname{Coh}(X_0)$ 对 $X_0$ 上完美复形的满子范畴
$\operatorname{Perf}(X_0)$ 的 dg-商范畴。
\end{definition}

另一个是 $W$ 的矩阵分解范畴 $\operatorname{Coh}(X,W)$：

\begin{definition}\label{def:matrix-factorization}
$W$ 在 $X$ 上的一个矩阵分解，由 $X$ 上一对相干层 $F_0,F_1$ 连同态射
$\phi_0\colon F_0\to F_1$ 与 $\phi_1\colon F_1\to F_0$ 组成，二者满足
$\phi_0\phi_1=W\cdot\id$ 与 $\phi_1\phi_0=W\cdot\id$。
\end{definition}

\transnote{易错预警：定义 \ref{def:sing-cat} 的「dg-商」指 dg 范畴意义下的 Drinfeld 商，不是普通范畴的
Verdier 商；它把完美复形整块压成零，剩下的非完美部分才携带奇点信息。}

% ===== 原书 p.88 =====
对该范畴中的一对对象 $E=(E_0,E_1,\phi^E)$ 与 $F=(F_0,F_1,\phi^F)$，态射复形由
\[
\Hom^0(E,F)=\Hom(E_0,F_0)\oplus\Hom(E_1,F_1),\qquad
\Hom^1(E,F)=\Hom(E_0,F_1)\oplus\Hom(E_1,F_0)
\]
给出，其微分为 $\de^i(f)=\phi^F\circ f-(-1)^if\circ\phi^E$。由 $X$ 上 $W$ 的矩阵分解
构成的 dg 导出范畴记作 $\operatorname{Coh}(X,W)$。

注意，$\operatorname{Sing}(X_0)$ 是 2-周期的 dg 范畴（因为 $\operatorname{Sing}$ 中每个
对象都有 2-周期的无穷解消），而 $\operatorname{Coh}(X,W)$ 是 $\C((\beta))$-线性的 dg 范畴。二者之间的
同构需要选一个由 $W$ 在 $\operatorname{Sing}(X_0)$ 上诱导的 2-周期性结构（不同的 $W$ 可以
有同一个零点轨迹 $X_0$）：在该 2-周期性自然变换处局部化，便得到微分的 $\Z/2$-分次范畴。
进一步的讨论见 [PS22]。

\transnote{含糊点拨：$\C((\beta))$ 是以 $\beta$ 为形式参数的 Laurent 幂级数域，$|\beta|=2$；
所谓 $\C((\beta))$-线性结构就是把「平移两次」当作可逆操作，因而分次只在模 2 意义下保留。}

具体地，该同构这样实现：取 $X_0$ 上的相干层 $F$，把它拉回到 $X$，得到局部自由层的两步
解消
\[
\begin{tikzcd}[column sep=2.3em, row sep=small]
0 \arrow[r] & E^{-1} \arrow[r, bend left=42, "\phi_{-1}"]
  & E^{0} \arrow[l, bend left=42, "\phi_0"] \arrow[r] & i^*F \arrow[r] & 0
\end{tikzcd}
\]
且 $\phi_{-1}\phi_0=W$、$\phi_0\phi_{-1}=W$：这给出相应的矩阵分解。拉回到 $X_0$，层
$E^i$ 仍是局部自由的；又因 $W$ 在 $X_0$ 上为零，得到正合列
\[
0\to F\to E^{-1}\to E^0\to F\to 0.
\]

\transnote{易错预警：解消图里的 $\phi_{-1}$ 与 $\phi_0$ 是下标，与 $E^{-1}$、$E^0$ 的上标
编号恰好错开一位；文字层会把上下标拍平，须按页图读作 $\phi_{-1}\colon E^{-1}\to E^0$ 与
$\phi_0\colon E^0\to E^{-1}$。}

% ===== 原书 p.89 =====
迭代此过程，便得到 $F$ 在 $X_0$ 上的无穷解消
\[
\cdots\to E^{-3}\to E^{-2}\to E^{-1}\to E^0\to F\to 0,
\]
它作为层的复形是 2-周期的。

可以通过给 $X$ 配一个使 $W$ 齐次的 $\mathbb G_m$ 作用，把这些范畴的 $\Z$-分次恢复出来。设 $G$ 是
作用在 $X$ 上的仿射代数群，$\chi$ 是 $G$ 的特征标，并使 $W$ 是 $\chi$-半不变的。

\begin{definition}\label{def:graded-mf}
分次矩阵分解范畴 $\operatorname{Coh}_G(X,W)$ 的对象，由 $X$ 上一对 $G$-等变相干层 $F_0,F_1$ 连同
$G$-不变态射 $\phi_1\colon F_1\to F_0$ 与 $\phi_0\colon F_0\to F_1(\chi)$ 组成，二者满足
$\phi_0\phi_1=W\cdot\id$ 与 $\phi_1(\chi)\phi_0=W\cdot\id$。两个对象之间的态射复形现由
\[
\Hom^{2n}(E,F)=\Hom(E_0,F_0(\chi^n))\oplus\Hom(E_1,F_1(\chi^n)),
\]
\[
\Hom^{2n+1}(E,F)=\Hom(E_0,F_1(\chi^n))\oplus\Hom(E_1,F_0(\chi^{n+1}))
\]
给出，微分与上面相同：$\de^i(f)=\phi^F\circ f-(-1)^if\circ\phi^E$。
\end{definition}

\transnote{含糊点拨：定义 \ref{def:graded-mf} 里两次复合的权不对称，$\phi_0\phi_1$ 不带 $\chi$，而
$\phi_1(\chi)\phi_0$ 带一个 $\chi$；这不是笔误，正是它使复合成为 $G$-不变的 $W\cdot\id$。}

对照 [Hir17, Definition 2.2]：那里把该范畴记作 $\operatorname{Coh}_G(X,\chi,W)$。本稿选择在记号中省略
特征标 $\chi$ 的选取；关于矩阵分解范畴及其间函子的其余约定，一律参看 [Hir17]。

眼下回到 $G=\mathbb G_m$ 的情形：此时该矩阵分解范畴是 $\Z$-分次的。在所得的三角范畴中，通常的
同调 $\Z/2$-分次将被延拓为一种「剪切」$\Z$-分次，其中 $[2]=(\chi)$。

\transnote{易错预警：「剪切」分次里的 $[2]=(\chi)$ 是约定，意思是每平移两次同时扭一个
特征标 $\chi$；不要把 $(\chi)$ 读成函数取值，它是等变线丛的记号。}

% ===== 原书 p.90 =====
\begin{definition}\label{def:equiv-sing}
等变奇点范畴 $\operatorname{Sing}_{\mathbb G_m}(X_0)$，定义为 $\operatorname{Coh}_{\mathbb G_m}(X_0)$（即 $X_0$ 上的
$\mathbb G_m$-等变相干层）对等变完美复形的满子范畴的 dg-商范畴。
\end{definition}

现在，一旦把来自 $\operatorname{Coh}_{\mathbb G_m}(X_0)$ 的解消上的等变结构计入，$\operatorname{Sing}(X_0)$
上的同调 $\Z$-分次就不再是 2-周期的：在前面对 2-周期解消的显式构造中可以看到，实际上
$E^\bullet[2]=E^\bullet(\chi)$。Orlov [Orl09a] 证明的 $\operatorname{Coh}(X,W)$ 与
$\operatorname{Sing}(X_0)$ 之间的等价，至此延拓为 dg 范畴的等价；下述定理由
[Hir17, Theorem 3.6] 取 $G=\mathbb G_m$ 得到。

\begin{theorem}[{$\Z$-分次 Orlov 定理}]\label{thm:zgraded-orlov}
给定 $X$ 上一个 $\mathbb G_m$ 作用，使 $W$ 是权 $1$ 的拟齐次函数，则存在 dg 范畴的等价
\[
\operatorname{Sing}_{\mathbb G_m}(X_0)\to\operatorname{Coh}_{\mathbb G_m}(X,W),
\]
它由包含 $X_0\hookrightarrow X$ 下的前推给出。
\end{theorem}

注意，由于 $W$ 在 $X_0$ 上为零，把 $X_0$ 上的层的复形前推，就得到 $X$ 上的矩阵分解。

\begin{lemma}\label{lem:6-1-6}
若 $Y$ 是任一带平凡 $\mathbb G_m$ 作用的光滑代数簇，则 $\operatorname{Coh}(Y)$ 与 $\operatorname{Coh}_{\mathbb G_m}(Y,0)$ 作为 dg
范畴等价。
\end{lemma}

见 [Hir17, Definition 2.3]。

\transnote{背景直觉：定理 \ref{thm:zgraded-orlov} 把「奇点范畴」换成「矩阵分解范畴」，代价是必须保留
$\mathbb G_m$-等变结构；等变结构记住的正是 $W$ 的权，从而把只在模 2 意义下存在的分次还原成
$\Z$-分次。}

% ===== 原书 p.91 =====
我们将用到的、Orlov [Orl06b] 的 Knörrer 周期性定理的基本 $\Z$-分次版本如下。

\begin{lemma}[{$\Z$-分次 Knörrer 周期性}]\label{lem:6-1-7}
若 $K_X$ 在纤维上带有权 $1$ 的 $\mathbb G_m$ 作用，则有 dg 范畴的等价
\[
\operatorname{Coh}(D)\xrightarrow{j_{D*}p_D^*}\operatorname{Sing}_{\mathbb G_m}(Z)
\xrightarrow{i_{Z*}}\operatorname{Coh}_{\mathbb G_m}(K_X,s(x)y),
\]
其中 $y$ 是 $K_X$ 的纤维上的坐标。
\end{lemma}

\begin{proof}
这由上述定理 \ref{thm:zgraded-orlov} 与 Hirano 对导出 Knörrer 周期性的 $\Z$-分次增强
[Hir17, Theorem 1.2] 合起来得到，其中取 $G=\mathbb G_m$、$\cE$ 为 $X$ 上的线丛 $\cK_X$、截面为
$s$。
\end{proof}

\transnote{易错预警：引理 \ref{lem:6-1-7} 的箭头方向容易读反。第一个函子把 $\operatorname{Coh}(D)$ 送到
$\operatorname{Sing}_{\mathbb G_m}(Z)$，第二个函子才送到 $\operatorname{Coh}_{\mathbb G_m}(K_X,s(x)y)$。
此外 $K_X$ 记典范丛的全空间，$\cK_X$ 记它作为层的版本，二者不可混用。}

我们还会用到它的若干升级版本。把 Hirano 文中编号 1.2 的定理用于 $G=\mathbb G_m$、$\cE$ 为 $K_X$ 上的
平凡线丛，便得到：

\begin{proposition}\label{prop:6-1-8}
存在 dg 范畴的等价
\[
j_{Z*}p_Z^*\colon\operatorname{Coh}_{\mathbb G_m}(Z,0)\to\operatorname{Coh}_{\mathbb G_m}(K_X\times\mathbb A^1,s(x)yz),
\]
这里：
\begin{itemize}
\item $\mathbb G_m$ 平凡地作用在 $K_X$ 上，并以权 $1$ 作用在 $\mathbb A^1$ 上；
\item $s$ 是 $\cK_X^{-1}$ 的截面，它定义环簇反典范除子；
\item $x$ 是 $X$ 上的坐标，$y$ 是 $K_X$ 的纤维上的坐标，$z$ 是 $\mathbb A^1$ 上的坐标；
\item $p_Z$ 是投影 $Z\times\mathbb A^1\to Z$，而 $j_Z$ 是包含
  $Z\times\mathbb A^1\to K_X\times\mathbb A^1$。
\end{itemize}
\end{proposition}

% ===== 原书 p.92 =====
同理，把 Hirano 定理用于 $X\times\mathbb A^1$ 上的 $\cE=\cK_X$、且 $\mathbb A^1$ 因子上带
权 $1$ 的 $\mathbb G_m$ 作用，便得到：

\begin{proposition}\label{prop:6-1-9}
存在 dg 范畴的等价
\[
j_{\tilde S*}p_{\tilde S}^*\colon\operatorname{Coh}_{\mathbb G_m}(\tilde S,0)
\to\operatorname{Coh}_{\mathbb G_m}(K_X\times\mathbb A^1,s(x)yz),
\]
其中 $\mathbb G_m$ 如前平凡地作用在 $K_X$ 上、并以权 $1$ 作用在 $\mathbb A^1$ 上；
$p_{\tilde S}$ 是投影 $\tilde S\times\mathbb A^1\to\tilde S$，而 $j_{\tilde S}$ 是包含
$K_X|_{\tilde S}\to K_X\times\mathbb A^1$。
\end{proposition}

这可由 Hirano 文中编号 1.2 的定理推出：为使 $zs(x)$ 给出 $\cK_X^{-1}$ 的一个 $\mathbb G_m$-不变截面，
线丛 $\cE=\cK_X$ 必须带有特征标 $\chi^{-1}$ 的等变结构；此时 $\cE\otimes\cO(\chi^1)$ 的
全空间的纤维在 $\mathbb G_m$ 作用下权为 $0$。

\transnote{含糊点拨：$Z$、$\tilde S$、$\hat Z$、$\tilde Z$ 这四个记号沿用第 1 章 §1.5 的
约定：$Z$ 是 $K_X$ 内的环簇反典范除子，$\tilde S$ 是 $X\times\mathbb A^1$ 内的环簇反典范
除子，$\hat Z$ 是 $K_X\times\mathbb A^1$ 内的环簇反典范除子，$\tilde Z$ 是粘合所得的导出
概型。$\hat Z$ 与 $\tilde Z$ 形状相近，极易读混。}

考虑把 $X\times\{0\}$ 沿 $D$ 粘到 $D\times\mathbb A^1[-1]$ 上所得的导出概型，我们把它记作
$\tilde Z$。

\begin{proposition}\label{prop:6-1-10}
有 dg 范畴的等价 $\operatorname{Coh}(\tilde Z)\cong\operatorname{Coh}_{\mathbb G_m}(\tilde S,0)$。
\end{proposition}

为给出证明，我们分若干步进行。

\begin{lemma}[{Koszul 对偶}]\label{lem:koszul-duality}
范畴 $\operatorname{Coh}(\mathbb A^1[-1])$ 等价于分次环 $\C[z]$ 上完美 dg-模范畴
$\operatorname{Perf}(\C[z])$，其中 $z$ 是上同调度为 $2$ 的变量。
\end{lemma}

\begin{proof}
例如见 [AG15, Corollary 5.1.10]，再取紧对象。
\end{proof}

% ===== 原书 p.93 =====
\begin{lemma}\label{lem:6-1-12}
范畴 $\operatorname{Coh}(\mathbb A^1[-1])$ 等价于分次矩阵分解范畴
$\operatorname{Coh}_{\mathbb G_m}(\mathbb A^1,0)$，其中 $\mathbb G_m$ 以权 $1$ 作用在
$\mathbb A^1$ 上。
\end{lemma}

\begin{proof}
这由上面 $\operatorname{Coh}(\mathbb A^1[-1])$ 与 $\operatorname{Perf}(\C[z])$ 的 Koszul 对偶等价，
与下述显式等价合起来得到：它把 dg-$\C[z]$-模 $(V^\bullet,d^\bullet)$ 送到矩阵分解
\[
F_0=\bigoplus_{n\in\Z}V^{2n}(\chi^{-n}),\qquad
F_1=\bigoplus_{m\in\Z}V^{2m+1}(\chi^{-m}),
\]
其中 $F_i$ 的 $\C[x]$-模结构由奇偶部分 $z\colon F_i\to F_i(\chi)$ 给出；而映射
$\varphi_0,\varphi_1$ 来自微分，$\varphi_0=\oplus_nd^{2n}(\chi^{-n})$、
$\varphi_1=\oplus_nd^{2n+1}(\chi^{-n})$。

类似地，若 $f^\bullet\colon(V^\bullet,d_V^\bullet)\to(W^\bullet,d_W^\bullet)$ 是
$\C[z]$-模的度 $2k$ 态射，则把 $f$ 限制到复形的奇、偶次数上，就得到 $\mathbb G_m$-等变
$\C[x]$-模的态射 $f_0\colon F_0(V)\to F_0(W)(\chi^k)$ 与
$f_1\colon F_1(V)\to F_1(W)(\chi^k)$，且它们与 $\varphi_V,\varphi_W$ 交换。但这恰是矩阵
分解 $F(V)$ 与 $F(W)$ 之间度 $2k$ 的态射（度 $2k+1$ 的情形同理）。只要 $(V^\bullet,d^\bullet)$
属于 $\operatorname{Perf}$，$F_0,F_1$ 就拟同构于相干的 $\C[x]$-模。不难验证，该函子保持
拟等价，并给出 dg 范畴的等价。
\end{proof}

\transnote{易错预警：从引理 \ref{lem:6-1-12} 起，原书把矩阵分解的映射记作 $\varphi_i$，而定义 \ref{def:matrix-factorization}、\ref{def:graded-mf} 记作 $\phi_i$；两类记号指同一对象，译本按页图分别保留，不强行统一。}

由此引理还可推出：

\begin{corollary}\label{cor:6-1-13}
$\operatorname{Coh}(D\times\mathbb A^1[-1])$ 与 $\operatorname{Coh}_{\mathbb G_m}(D\times\mathbb A^1,0)$ 之间有 dg 范畴的
等价，其中 $\mathbb G_m$ 平凡地作用在 $D$ 因子上，并以权 $1$ 作用在 $\mathbb A^1$ 上。
\end{corollary}

下述结果也可视为引理 \ref{lem:6-1-12} 的一个更一般的形式：

% ===== 原书 p.94 =====
\begin{proposition}\label{prop:6-1-14}
对带 $G$ 作用的概型 $Y$ 与特征标 $\chi$ 的选取，矩阵分解范畴 $\operatorname{Coh}_G(Y,0)$ 与等变相干层
范畴 $\operatorname{Coh}_G(Y\times\mathbb A^1[-1])$ 之间有 dg 范畴的等价，其中 $G$ 经 $\chi$ 作用在
$\mathbb A^1$ 上。
\end{proposition}

\begin{proof}
$\operatorname{Coh}_G(Y\times\mathbb A^1[-1])$ 的一个对象，由 $Y$ 上的 $G$-等变相干层复形 $\cF^\bullet$
连同度 $(-1)$ 的 $G$-等变映射 $t\colon\cF^\bullet\to\cF^{\bullet-1}(\chi)$ 组成，并满足
$t^2=0$：这就是 $\mathbb A^1[-1]=\operatorname{Spec}\C[t]$（$|t|=-1$）的权 $\chi$ 坐标
$t$ 在 $\cF$ 上的作用。给定这样的对象，我们构造 $0$ 的一个矩阵分解，取
\[
F_0=\bigoplus_{n\in\Z}\cF^{2n}(\chi^{-n}),\qquad
F_1=\bigoplus_{m\in\Z}\cF^{2m+1}(\chi^{-m}),
\]
并有 $\varphi_1=t\colon F_1\to F_0$、$\varphi_0=t\colon F_0\to F_1(\chi)$。类似地，度
$2k$ 的态射 $f^\bullet\colon\cF^\bullet\to\cG^\bullet$ 通过
$f_0\colon F_0\to G_0(\chi^k)$ 与 $f_1\colon F_1\to G_1(\chi^k)$ 给出 $Y$ 上 $G$-等变
相干层的态射，且与 $\varphi_F,\varphi_G$ 交换，从而给出 $\operatorname{Coh}_G(Y,0)$ 中度 $2k$ 的态射
（度 $2k+1$ 的情形同理）。如上所述，不难验证这给出 dg 范畴的等价。
\end{proof}

\begin{proposition}\label{prop:6-1-15}
函子 $\operatorname{Coh}_G(\cdot,0)$ 关于带 $G$ 作用的概型的 proper 等变态射满足下降。
\end{proposition}

\begin{proof}
由于命题 \ref{prop:6-1-14}，可以先应用 [GR17, Proposition 4.7.2.2] 的下降结果，再取紧对象。该结果
虽是对 aft 概型陈述的，但我们需要对叠的陈述由下降得到。
\end{proof}

最后，我们可以证明命题 \ref{prop:6-1-10}：

% ===== 原书 p.95 =====
\begin{proof}[命题 \ref{prop:6-1-10} 的证明]
我们将证明，$\operatorname{Coh}_{\mathbb G_m}(\tilde S,0)$ 可以写成给出 $\operatorname{Coh}(\tilde Z)$ 的同一个图的推出
（见命题 \ref{prop:7-2-3}）。由引理 \ref{lem:6-1-6} 与推论 \ref{cor:6-1-13} 可知，可以把推出图中的三个范畴分别写成矩阵
分解范畴：
\begin{itemize}
\item $\operatorname{Coh}(D)\simeq\operatorname{Coh}_{\mathbb G_m}(D,0)$；
\item $\operatorname{Coh}(D\times\mathbb A^1[-1])\simeq\operatorname{Coh}_{\mathbb G_m}(D\times\mathbb A^1,0)$；
\item $\operatorname{Coh}(X)\simeq\operatorname{Coh}_{\mathbb G_m}(X,0)$；
\end{itemize}
图中相应的函子由矩阵分解上的前推函子给出。于是由命题 \ref{prop:6-1-15}，
\[
\begin{tikzcd}[column sep=3.4em, row sep=2.2em]
\operatorname{Coh}_{\mathbb G_m}(D,0) \arrow[r] \arrow[d] & \operatorname{Coh}_{\mathbb G_m}(X,0) \arrow[d] \\
\operatorname{Coh}_{\mathbb G_m}(D\times\mathbb A^1,0) \arrow[r] & \operatorname{Coh}_{\mathbb G_m}(\tilde S,0)
\end{tikzcd}
\]
是推出方块，其中的箭头由前推给出。
\end{proof}

\transnote{易错预警：该推出方块的四个角是 $\operatorname{Coh}_{\mathbb G_m}(\cdot,0)$，不是普通的 $\operatorname{Coh}$；文字
层只留下四个名字，箭头与方块须按页图读。方块中上行的两个范畴来自 $D$ 与 $X$，下行的两个
来自 $D\times\mathbb A^1$ 与 $\tilde S$，切勿把 $\tilde S$ 换成 $\hat Z$。}

结合命题 \ref{prop:6-1-8}、\ref{prop:6-1-9}、\ref{prop:6-1-10}，我们得到：

\begin{theorem}[{Knörrer 周期性}]\label{thm:6-1-16}
设 $\mathbb G_m$ 平凡地作用在 $K_X$ 上，并以权 $1$ 作用在 $K_X\times\mathbb A^1$ 的 $\mathbb A^1$
纤维上，则有 dg 范畴的等价：
\[
\operatorname{Coh}(Z)\longrightarrow\operatorname{Coh}_{\mathbb G_m}(K_X\times\mathbb A^1,s(x)yz)
\longleftarrow\operatorname{Coh}(\tilde Z).
\]
\end{theorem}

\transnote{背景直觉：定理 \ref{thm:6-1-16} 是说 $\operatorname{Coh}(Z)$ 与 $\operatorname{Coh}(\tilde Z)$ 并不相等，但二者经由
同一个中间范畴「稳定化」后一致；这正是第 1 章所说的公共 Knörrer 稳定化。原书此处的箭头
未标函子名，按前三个命题即 $j_{Z*}p_Z^*$ 与 $j_{\tilde S*}p_{\tilde S}^*$。}

% ===== 原书 p.96 =====
\subsection{双分次}

本节考虑把上述结果延拓到分次 dg 范畴：即对分次向量空间的范畴 $\operatorname{Coh}(B\mathbb G_m)$ 线性的范畴。
我们把分次 dg 范畴之间保持分次的函子称作分次 dg 函子。

设考虑 $K_X\times\mathbb A^1$ 上的 $\mathbb G_m\times\mathbb G_m$ 作用：第一个 $\mathbb G_m$ 以权 $1$ 作用在
$\mathbb A^1$ 因子上（平凡地作用在 $K_X$ 上），第二个 $\mathbb G_m$ 以权 $1$ 作用在 $K_X$ 的纤维
上（平凡地作用在 $X\times\mathbb A^1$ 上）。我们让 $Z$、$\tilde S$ 带有诱导的
$\mathbb G_m\times\mathbb G_m$ 作用。第一个 $\mathbb G_m$ 因子作用的权给出同调分次，第二个因子的权给出辅助分次。

\transnote{易错预警：原文此处印作「the second second factor」，第二个 second 是重复排印，
应读作「第二个因子」；页图亦如此，非 OCR 误读。}

\begin{theorem}\label{thm:6-2-1}
存在分次 dg 范畴的等价：
\[
\operatorname{Coh}_{\mathbb G_m\times\mathbb G_m}(Z,0)\xrightarrow{j_{Z*}p_Z^*}
\operatorname{Coh}_{\mathbb G_m\times\mathbb G_m}(K_X\times\mathbb A^1,s(x)yz)
\xleftarrow{j_{\tilde S*}p_{\tilde S}^*}\operatorname{Coh}_{\mathbb G_m\times\mathbb G_m}(\tilde S,0),
\]
它们由上述同一批 Knörrer 周期性函子诱导。
\end{theorem}

\begin{proof}
由 [Hir17, Theorem 1.2] 取 $G=\mathbb G_m\times\mathbb G_m$ 即得。
\end{proof}

下一条结果同样由命题 \ref{prop:6-1-15} 取 $G=\mathbb G_m\times\mathbb G_m$ 得到。

\begin{proposition}\label{prop:6-2-2}
有分次 dg 范畴的推出方块：
\[
\begin{tikzcd}[column sep=3.4em, row sep=2.2em]
\operatorname{Coh}_{\mathbb G_m^2}(D,0) \arrow[r] \arrow[d] & \operatorname{Coh}_{\mathbb G_m^2}(X,0) \arrow[d] \\
\operatorname{Coh}_{\mathbb G_m^2}(D\times\mathbb A^1,0) \arrow[r] & \operatorname{Coh}_{\mathbb G_m^2}(\tilde S,0)
\end{tikzcd}
\]
% ===== 原书 p.97 =====
其中第二个 $\mathbb G_m$ 因子平凡地作用，第一个因子以权 $1$ 作用在 $\mathbb A^1$ 因子上。
\end{proposition}

我们还需要引理 \ref{lem:6-1-6} 的一个升级形式：

\begin{lemma}\label{lem:6-2-3}
若 $Y$ 是任一带 $\mathbb G_m$ 作用的光滑代数簇，则等变相干层范畴 $\operatorname{Coh}_{\mathbb G_m}(Y)$ 与
$\operatorname{Coh}_{\mathbb G_m^2}(Y,0)$（第二个 $\mathbb G_m$ 因子平凡地作用）作为分次 dg 范畴等价。
\end{lemma}
